% ---
% id: reductions/lem:latsum
% kind: lem
% title: "Partial sums on the declared lattice are exact"
% summary: >-
%   If κ + ⌈log₂ L⌉ ≤ 53, every partial sum of L entries on the declared
%   lattice is a binary64 number and every addition forming it is exact. A
%   family on the lattice over such an alphabet therefore meets the
%   rounds-nowhere condition by construction.
% prominence: technical
% status: accepted
% uses: [def:exactadm]
% mentions: [foundations/lem:exact-accumulation]
% assumes: []
% refs: []
% provenance: newly-derived
% lean:
% ---

A sum \tref{rounds nowhere} exactly when every partial sum is a binary64
number (\xref{foundations/lem:exact-accumulation}). On the declared lattice
that holds as long as the alphabet is small enough for the partial sums to
fit in the significand.

\begin{lemma}\label{lem:latsum}
Let $x_1,\dots,x_L$ satisfy \eqref{eq:latinv} and let
$\Sigma_j\coloneqq\sum_{i\le j}x_i$, $1\le j\le L$, be the partial sums of the
accumulation in that order. If
\begin{equation}\label{eq:latsumcond}
  \kappa+\big\lceil\log_2L\big\rceil\ \le\ 53 ,
\end{equation}
then every $\Sigma_j$ is exactly a binary64 number and every one of the $L-1$
additions forming it is exact in binary64. Consequently a family whose entries
satisfy \eqref{eq:latinv}, over an alphabet satisfying
\eqref{eq:latsumcond}, meets condition~(e) of Definition~\ref{def:exactadm} by
construction.
\end{lemma}

\begin{proof}
Work one real component of one entry position at a time; the accumulation is
entrywise and the real and imaginary parts are accumulated independently. Each
$x_i$ is an integer multiple of $2^{-\kappa}$, so each $\Sigma_j$ is one too,
and $|\Sigma_j|\le\sum_{i\le j}|x_i|\le L$. Hence $2^{\kappa}\Sigma_j$ is an
integer with $\big|2^{\kappa}\Sigma_j\big|\le L\,2^{\kappa}\le2^{53}$ by
\eqref{eq:latsumcond}. Every integer of modulus at most $2^{53}$ is exactly
representable in binary64, and multiplication by the power of two
$2^{-\kappa}$ is exact and changes no significand, so $\Sigma_j$ is a binary64
number; and $2^{-\kappa}\ge2^{-53}$ with $|\Sigma_j|\le L$ keeps it far
inside the normal range, so no subnormal or overflow case arises. An IEEE-754
addition returns the correctly rounded value of the exact sum of its operands;
here that exact sum is $\Sigma_j$, which is representable, so the rounding is
the identity and the addition is exact.
\end{proof}
